% Section 6.7, translated from sv/chapters/06/exercises.tex.

\section{Exercises}
\label{c6:sec:uppgifter}

\begin{aside}
\textbf{How to work with the exercises.} Parts~1 and~2 are done during the lesson: first on your
own, a few minutes per exercise, then together. Part~3 is extra exercises for further practice
after the lesson. Most of them can be done in your head or with pencil and paper.

\facitnotis{L06}
\end{aside}

% ------------------------------------------------------------------------------------------
\uppgiftsdel{Part 1}{Review exercises}

\uppgift{1.1}{Laws of exponents}
Simplify, with positive exponents in the answer.
\begin{delar}(3)
  \task $\dfrac{10^6 \cdot 10^{-3}}{10^2}$
  \task $(2 \cdot 10^3)^2$
  \task $\sqrt{4 \times 10^6}$
\end{delar}

\uppgift{1.2}{Scientific notation in a circuit calculation}
A circuit has $U = 3.3\,\text{V}$ and $R = 4.7\,\text{k}\Omega = 4.7 \times 10^3\,\Omega$.
\begin{parts}
  \item Calculate the current $I = \dfrac{U}{R}$ and give the answer in mA in scientific
    notation.
  \item Calculate the power $P = \dfrac{U^2}{R}$ and give the answer in mW.
\end{parts}

\uppgift{1.3}{Square root}
\begin{parts}
  \item Calculate $\sqrt{144}$.
  \item Calculate $\sqrt{0.0049}$.
  \item A resistor dissipates $P = 2\,\text{W}$ at the voltage $U$, and $R = 200\,\Omega$.
    Calculate $U$ using $P = \dfrac{U^2}{R}$.
\end{parts}

% ------------------------------------------------------------------------------------------
\uppgiftsdel{Part 2}{New material}

\uppgift{2.1}{A pure quadratic equation}
The current $I$ (in amperes) satisfies
\[
  50I^2 = 8.
\]
\begin{parts}
  \item Solve the equation for $I$.
  \item Explain why the negative root should be rejected in this electrical context.
\end{parts}

\uppgift{2.2}{Quadratic equations with the PQ-formula}
Solve the following equations with the PQ-formula.
\begin{delar}(3)
  \task $x^2 - 5x + 6 = 0$
  \task $x^2 + 3x - 10 = 0$
  \task $x^2 - 6x + 9 = 0$
\end{delar}

\uppgift{2.3}{Quadratic equations with the ABC-formula}
Solve with the ABC-formula.
\begin{delar}(3)
  \task $2x^2 - 7x + 3 = 0$
  \task $3x^2 + 5x - 2 = 0$
\end{delar}

\uppgift{2.4}{Resistors in series}
Two resistors $R_1$ and $R_2$ in series satisfy
\[
  R_1 + R_2 = 13\,\Omega \qquad \text{and} \qquad R_1 \cdot R_2 = 36\,\Omega^2.
\]
\begin{parts}
  \item Use Vieta's formulas to set up a quadratic equation with $R_1$ as the unknown.
  \item Solve the equation and give both resistances.
\end{parts}

% ------------------------------------------------------------------------------------------
\uppgiftsdel{Part 3}{Extra exercises}
The exercises below are extra practice, best done on your own after the lesson.

\uppgift{3.1}{Pure quadratic equations}
Solve the equations.
\begin{delar}(3)
  \task $x^2 = 49$
  \task $3x^2 = 27$
  \task $x^2 - 5 = 20$
  \task $2x^2 + 8 = 0$
  \task $\dfrac{x^2}{4} = 9$
\end{delar}

\uppgift{3.2}{The factorisation method}
Solve by factorising and using the zero-product rule.
\begin{delar}(3)
  \task $x^2 - 7x = 0$
  \task $x^2 - 16 = 0$
  \task $x^2 + 5x + 6 = 0$
  \task $2x^2 - 8x = 0$
  \task $x^2 - 4x + 4 = 0$
\end{delar}

\uppgift{3.3}{The PQ-formula}
Solve with the PQ-formula.
\begin{delar}(2)
  \task $x^2 + 2x - 8 = 0$
  \task $x^2 - 8x + 12 = 0$
  \task $x^2 + x - 12 = 0$
  \task $x^2 - 2x - 15 = 0$
\end{delar}

\uppgift{3.4}{Normalise first}
The equations below are not normalised. First divide so that $a = 1$, then solve with the
PQ-formula.
\begin{delar}(3)
  \task $2x^2 - 10x + 12 = 0$
  \task $3x^2 + 12x - 15 = 0$
  \task $-x^2 + 4x - 3 = 0$
\end{delar}

\uppgift{3.5}{The discriminant}
Calculate the discriminant $D = b^2 - 4ac$ and give the number of real roots \textbf{without}
solving the equation.
\begin{delar}(3)
  \task $x^2 - 6x + 5 = 0$
  \task $x^2 + 4x + 4 = 0$
  \task $x^2 + x + 3 = 0$
  \task $2x^2 - 4x + 2 = 0$
  \task $3x^2 + 2x - 1 = 0$
\end{delar}

\uppgift{3.6}{Completing the square}
Solve by completing the square.
\begin{delar}(3)
  \task $x^2 + 4x - 5 = 0$
  \task $x^2 - 10x + 21 = 0$
  \task $x^2 + 6x + 4 = 0$
\end{delar}

\uppgift{3.7}{Vieta's formulas}
\begin{parts}
  \item A normalised quadratic equation has the roots $x_1 = 4$ and $x_2 = -2$. Find $p$ and $q$
    and write down the equation.
  \item Check your answer to a) with the PQ-formula.
  \item Use Vieta's formulas to check quickly that $x_1 = 4$ and $x_2 = 5$ are roots of
    $x^2 - 9x + 20 = 0$.
  \item Can a normalised quadratic equation have the double root $x = 3$? If so, write it down.
\end{parts}

\uppgift{3.8}{Current from power}
The power in a resistor is given by $P = RI^2$.
\begin{parts}
  \item $P = 12.5\,\text{W}$ and $R = 50\,\Omega$. Calculate $I$.
  \item $P = 0.18\,\text{W}$ and $R = 200\,\Omega$. Calculate $I$ in mA.
  \item A resistor is marked $470\,\Omega$ and $0.5\,\text{W}$. Calculate the largest permitted
    current.
  \item Why is the negative root rejected in every part?
\end{parts}

\uppgift{3.9}{Voltage from power}
The power can also be written $P = \dfrac{U^2}{R}$.
\begin{parts}
  \item $P = 5\,\text{W}$ and $R = 20\,\Omega$. Calculate $U$.
  \item $P = 0.125\,\text{W}$ and $R = 2\,000\,\Omega$. Calculate $U$.
  \item A heating element is to give off $P = 1\,000\,\text{W}$ at $U = 230\,\text{V}$. What
    resistance is needed?
\end{parts}

\uppgift{3.10}{Two resistors from their sum and product}
\begin{parts}
  \item Two resistors satisfy $R_1 + R_2 = 20\,\Omega$ and $R_1 R_2 = 96\,\Omega^2$. Find the
    resistances.
  \item Two resistors satisfy $R_1 + R_2 = 15\,\Omega$ and $R_1 R_2 = 56\,\Omega^2$. Find the
    resistances.
  \item Can two resistors have the sum $10\,\Omega$ and the product $30\,\Omega^2$? Justify your
    answer with the discriminant.
\end{parts}

\uppgift{3.11}{Series and parallel resistance at once}
Two resistors together give $10\,\Omega$ in series and $2.4\,\Omega$ in parallel.
\begin{parts}
  \item Show that $R_1 R_2 = 24\,\Omega^2$.
  \item Set up and solve the quadratic equation for the resistances.
  \item Check both the series and the parallel resistance.
\end{parts}

\uppgift{3.12}{Power in a load}
An unknown resistor $R$ is connected in series with $R_1 = 4\,\Omega$ and supplied with
$E = 12\,\text{V}$. The power developed in $R$ is to be $P = 8\,\text{W}$. The current in the
circuit is $I = \dfrac{E}{R_1 + R}$ and the power in $R$ is $P = RI^2$.
\begin{parts}
  \item Show that the requirement leads to the equation $R^2 - 10R + 16 = 0$.
  \item Solve the equation.
  \item Check \textbf{both} solutions by calculating the current and the power.
  \item Calculate the power in $R$ when $R = 4\,\Omega$ and compare with the two solutions.
\end{parts}

% Kept whole on one page: its last part was left alone at the top of the next.
\needspace{16\baselineskip}
\uppgift{3.13}{True or false}
Decide whether the statement is true or false and justify your answer briefly.
\begin{parts}
  \item Every quadratic equation has two real roots.
  \item The equation $x^2 = 16$ has the solution $x = 4$.
  \item The PQ-formula can be used directly on $2x^2 + 4x - 6 = 0$.
  \item If $D = 0$, the two roots coincide.
  \item If a product of two factors is zero, at least one of the factors must be zero.
  \item The equation $x^2 + 4 = 0$ has no real solutions.
\end{parts}
